\documentclass{article}
\usepackage{graphicx} % Required for inserting images
\usepackage[T1]{fontenc}
\usepackage{amsmath}
\usepackage[margin=1in]{geometry}
\usepackage{booktabs}
\usepackage{tabularx}
\usepackage[hidelinks]{hyperref}

\title{Lesson 2: Options from First Principles}
\author{Analyst onboarding \textperiodcentered{} Lesson 2 of 10 \textperiodcentered{} L/S Gamma Desk, QUANTT}
\date{2026-10-02}

\begin{document}

\maketitle

\textbf{Goal:} know exactly what an option is, what it pays, and why our desk trades them.

\textbf{You need:} Lesson 1 (SPY, long/\allowbreak{}short, bid/\allowbreak{}ask, volatility).

\section{What an option is}

An option is a contract that gives its buyer a \textbf{right, but not an obligation}, to buy or sell something at a price agreed today, up to a set date. The buyer pays the seller for that right up front. This payment is the \textbf{premium}.

Think of it like a concert-ticket reservation: you pay \$20 now to lock in the right to buy a \$100 ticket next month. If prices jump to \$300, your right is valuable. If they fall to \$50, you walk away and lose only the \$20.

\begin{itemize}

\item \textbf{Call option:} the right to \textbf{buy} at the agreed price.

\item \textbf{Put option:} the right to \textbf{sell} at the agreed price.

\item \textbf{Strike price (K):} the agreed price.

\item \textbf{Expiry:} the last day the right exists. \textbf{DTE} means \emph{days to expiry}.

\item \textbf{Exercise:} using the right, i.e. actually buying or selling at the strike.

\item \textbf{Underlying:} what the option is on. For us, \textbf{SPY}.

\end{itemize}

\textbf{$\Rightarrow$} The buyer risks only the premium and keeps the upside. The seller collects the premium and takes on the risk.

\section{The details of SPY options}

\begin{itemize}

\item \textbf{One contract covers 100 shares.} A quoted price of \$13.40 costs \$13.40 \ensuremath{\times} 100 = \$1,340.

\item \textbf{American style:} they can be exercised on any trading day up to expiry.

\item \textbf{They trade with a bid and an ask} like SPY itself, and the spread is a cost.

\item \textbf{There are expiries for every trading day.} Our desk uses Friday expiries about one month out.

\end{itemize}

\section{What an option pays at expiry}

On the expiry day the option is worth only what exercising it would pay:

\[\text{Call value} = \max(S - K,\ 0) \qquad \text{Put value} = \max(K - S,\ 0)\]

Here $S$ is SPY's price at expiry and ``max'' means ``take the larger of the two''. If exercising would lose money, you simply don't, so the value is 0. \textbf{Profit = value \ensuremath{-} premium paid.}

\textbf{Example.} SPY trades at 760. The 761-strike call expiring in 30 days costs \$13.40, and the 761 put costs \$11.20 (real fills from our paper account).

\begin{center}\small
\begin{tabularx}{\linewidth}{>{\raggedright\arraybackslash}X >{\raggedright\arraybackslash}X >{\raggedright\arraybackslash}X l}
\toprule
\textbf{SPY at expiry} & \textbf{Bought the call} & \textbf{Bought the put} & \textbf{Sold the put} \\
\midrule
740 & \ensuremath{-}\$1,340 & +\$980 & \ensuremath{-}\$980 \\
761 & \ensuremath{-}\$1,340 & \ensuremath{-}\$1,120 & +\$1,120 \\
780 & +\$560 & \ensuremath{-}\$1,120 & +\$1,120 \\
\bottomrule
\end{tabularx}
\end{center}

\begin{itemize}

\item \textbf{Bought the call:} the most it can lose is \$1,340. It breaks even if SPY ends at 761 + 13.40 = 774.40.

\item \textbf{Bought the put:} the most it can lose is \$1,120. It breaks even at 761 \ensuremath{-} 11.20 = 749.80.

\item \textbf{Sold the put:} it keeps the whole \$1,120 if SPY ends above 761, and loses more the further SPY falls.

\end{itemize}

\textbf{$\Rightarrow$} Buying options is a bet on a \textbf{big} move. Selling options is a bet on a \textbf{small} move. This one idea drives everything our desk does.

\section{Where the strike sits: moneyness}

\begin{itemize}

\item \textbf{At the money (ATM):} the strike is near SPY's current price. Our desk trades ATM options.

\item \textbf{In the money (ITM):} exercising now would pay something (a call with K below SPY, a put with K above).

\item \textbf{Out of the money (OTM):} exercising now would pay nothing (a call with K above SPY, a put with K below).

\end{itemize}

\section{Why an option is worth more than its payoff today}

Before expiry, an option's price has two parts:

\begin{itemize}

\item \textbf{Intrinsic value:} what it would pay if exercised right now. The 761 call with SPY at 760 has none.

\item \textbf{Time value:} everything else. The call still costs \$13.40 because SPY might rise a lot in 30 days.

\end{itemize}

Time value melts to zero by expiry. How much the market charges for it depends on \textbf{how much it expects SPY to move}, i.e. on volatility. Lesson 3 turns that into a number.

\section{Buyers and sellers face very different risks}

\begin{center}\small
\begin{tabularx}{\linewidth}{>{\raggedright\arraybackslash}X >{\raggedright\arraybackslash}X >{\raggedright\arraybackslash}X}
\toprule
\textbf{} & \textbf{Buyer (long the option)} & \textbf{Seller (short the option)} \\
\midrule
Premium & Pays it & Receives it \\
Worst case & Loses the premium & Large losses (unlimited when selling calls) \\
Hopes for & Big moves & Small moves \\
Each day that passes & Costs money & Earns money \\
\bottomrule
\end{tabularx}
\end{center}

\textbf{$\Rightarrow$} Sellers earn small, steady income and sometimes take a big loss. Buyers pay a small, steady cost and sometimes win big. Choosing \textbf{when to be which} is our job.

\section{Why use options at all?}

\begin{enumerate}

\item \textbf{Curved payoffs:} gains grow faster as the move grows, while losses are capped at the premium. Shares pay in a straight line.

\item \textbf{A bet on movement itself:} options let you bet on \emph{how much} SPY moves, not \emph{which way}. Lesson 5 shows how we remove the direction completely.

\item \textbf{Leverage:} \$1,340 gives exposure to 100 shares worth about \$76,000. That cuts both ways.

\end{enumerate}

Our desk uses options for reason 2: we are \textbf{volatility traders}, not directional traders.

\section{Words to know}

\begin{itemize}

\item \textbf{Call /\allowbreak{} put:} right to buy /\allowbreak{} sell. \textbf{Strike:} the agreed price. \textbf{Premium:} the option's price.

\item \textbf{Expiry, DTE:} last day, days left. \textbf{Exercise:} use the right.

\item \textbf{ATM /\allowbreak{} ITM /\allowbreak{} OTM:} strike near /\allowbreak{} in favour of /\allowbreak{} against the current price.

\item \textbf{Intrinsic /\allowbreak{} time value:} payoff now /\allowbreak{} the extra the market charges for possible movement.

\end{itemize}

\section{Check yourself}

\begin{enumerate}

\item You buy the 761 put for \$11.20 and SPY finishes at 755. What do you make or lose on one contract?

\item Why is the 761 call worth \$13.40 when exercising it today pays nothing?

\item If SPY doesn't move at all for 30 days, who makes money: the call buyer or the put seller?

\end{enumerate}

\emph{Answers: (1) (761 \ensuremath{-} 755 \ensuremath{-} 11.20) \ensuremath{\times} 100 = \ensuremath{-}\$520. (2) Time value: SPY might rise a lot before expiry. (3) The put seller keeps the premium; the call buyer loses it.}

\end{document}
